As medidas são calculadas no texto normalizado comum. Denotam-se por $C$ os caracteres não brancos, por $L$ as letras Unicode, por $T$ os segmentos lexicais, por $U$ os tipos distintos e por $B$ as linhas não vazias. Segmentos são sequências de letras que admitem apóstrofo ou hífen interno, após uniformização de caixa por \textit{casefold}. Razões com denominador nulo recebem zero. O indicador 21 registra a ausência de segmentos.

\begingroup\footnotesize\singlespacing
\begin{longtable}{rP{.31\linewidth}P{.56\linewidth}}
\caption{Definições operacionais dos 21 descritores de letras.}\label{tab:descritores}\\
\toprule & Descritor & Definição \\\midrule\endfirsthead
\toprule & Descritor & Definição (continuação) \\\midrule\endhead
\bottomrule\endfoot
\bottomrule\endlastfoot
1 & Caracteres totais & Pontos de código, incluindo espaços e quebras internas.\\
2 & Caracteres não brancos & Contagem $C$.\\
3 & Linhas não vazias & Contagem $B$.\\
4 & Estrofes operacionais & Blocos não vazios separados por uma ou mais linhas vazias.\\
5 & Segmentos lexicais & Contagem $T$.\\
6 & Tipos lexicais & Contagem $U$, após casefold.\\
7 & Comprimento médio de segmento & Média de caracteres dos segmentos após casefold.\\
8 & Dispersão de comprimento de segmento & Desvio com divisor $T$.\\
9 & Razão tipo/segmento & $U/T$.\\
10 & Razão de tipos únicos em frequência & Número de tipos com uma ocorrência dividido por $U$.\\
11 & Entropia de segmentos & $-\sum_{t=1}^{U}p_t\log_2p_t$, com frequências relativas $p_t$.\\
12 & Comprimento médio de linha & Média de caracteres não brancos nas linhas não vazias.\\
13 & Dispersão de comprimento de linha & Desvio com divisor $B$.\\
14 & Fração de linhas repetidas & $(B-\text{número de linhas distintas})/B$, após casefold e uniformização de espaços.\\
15 & Fração de letras & $L/C$, com categoria Unicode L.\\
16 & Fração de dígitos & Dígitos decimais Unicode divididos por $C$.\\
17 & Fração de pontuação & Caracteres de categoria Unicode P divididos por $C$.\\
18 & Fração de maiúsculas & Letras maiúsculas divididas por $L$.\\
19 & Linhas com interrogação & Fração das $B$ linhas com interrogação simples ou de largura completa.\\
20 & Linhas com exclamação & Fração das $B$ linhas com exclamação simples ou de largura completa.\\
21 & Ausência de segmentos & Indicador de $T=0$.\\
\end{longtable}
\endgroup

A entropia segue a definição discreta em bits de \citeonline{cover2006}, aplicada aqui aos segmentos. Ela e a razão tipo/segmento dependem da extensão e da segmentação. O indicador 21 permite registrar a ausência de segmentos na regra adotada.
